\chapter{Analyse et discussion}
\label{chap:chap_4}

La principale contrainte du système réside dans la visibilité limitée au niveau du chiffrement, notamment pour les protocoles sécurisés tels que SSH ou HTTPS, dont le contenu reste inaccessible sans inspection TLS.  
De plus, le déploiement sur une unique machine virtuelle limite la scalabilité et ne permet pas de simuler un environnement réseau distribué réel.  
Enfin, l’absence d’une corrélation automatique inter-scénarios ou d’un module d’analyse comportementale réduit la capacité de détection des attaques multi-étapes.

\bigskip
\noindent
Améliorations possibles :  
Une évolution naturelle du projet consisterait à intégrer des outils de réponse et d’automatisation tels que TheHive ou Cortex afin de créer un mini-SOAR (Security Orchestration, Automation and Response).  
L’ajout d’un module de machine learning pour la détection comportementale pourrait également améliorer la précision des alertes et réduire les faux positifs.  
Enfin, la mise en place d’un cluster Elasticsearch distribué et d’un Wazuh Manager offrirait une meilleure résilience et une supervision multi-nœuds.

\bigskip
\noindent
Perspectives et veille technologique :  
Ce projet constitue une base solide pour explorer des approches modernes telles que :
\begin{itemize}
    \item L’intégration d’outils de threat intelligence (MISP, OpenCTI) pour enrichir les alertes avec des indicateurs de compromission connus (IoC) ;
    \item L’adoption d’un pipeline de logs unifié avec Logstash ou Fluentd pour améliorer la normalisation et la performance d’ingestion ;
    \item L’utilisation d’outils de visualisation avancée (Grafana, OpenSearch Dashboards) pour compléter l’analyse de Kibana.
\end{itemize}
